\chapter{Протидія злоякісному коду у спеціалізованих середовищах}
\chaptershort{Спеціалізовані середовища}
\label{chap:ch14}

Методи й засоби попередніх розділів орієнтовані переважно на класичні обчислювальні системи (робочі станції, сервери загального призначення). Проте ЗПК дедалі частіше націлюється на \emph{спеціалізовані} середовища --- промислові системи керування, Інтернет речей, мобільні платформи, хмару, --- кожне з яких має власну модель загроз, обмеження й вектори, через що прямий перенос «настільних» підходів часто неможливий. Для України це питання має критичне значення: атаки на промислову автоматику енергетики (Industroyer) і руйнівні кампанії проти критичної інфраструктури показали, що спеціалізовані середовища є пріоритетними цілями. Розділ розглядає спільні принципи протидії в спеціалізованих середовищах, захист промислових систем (ICS/SCADA), безпеку Інтернету речей, мобільних платформ і хмарних та контейнерних середовищ, і завершується кейсом захисту промислового об'єкта. Виклад спирається на моделювання загроз (розділ~\ref{chap:ch04}), виявлення (розділи~\ref{chap:ch06}--\ref{chap:ch11}) і розвідку загроз (розділ~\ref{chap:ch12}), адаптуючи їх до специфіки кожного середовища.

\section{Спільні принципи й виклики}\label{sec:ch14-common}

Попри відмінності, спеціалізовані середовища поділяють спільні виклики, що відрізняють їх від класичних систем (табл.~\thechapter.1): \emph{обмежені ресурси} (неможливість встановити повноцінний агент захисту); \emph{специфічні ОС і протоколи} (сигнатури й засоби для класичних систем незастосовні); \emph{довгий життєвий цикл і складність оновлення} (пристрої працюють роками без патчів); \emph{висока ціна збою} (у промисловості й медицині --- фізична безпека людей); \emph{масовість і слабка видимість} (мільярди пристроїв IoT поза моніторингом). Ці чинники зміщують акцент із виявлення на кінцевій точці (часто неможливого) на мережевий рівень, сегментацію й профілактику.

\tabhead{Таблиця~\thechapter.1 --- Виклики протидії ЗПК у спеціалізованих середовищах}

\begin{center}\footnotesize
\begin{tabularx}{\linewidth}{L{22mm} Y L{22mm}}
\toprule
Виклик & Наслідок & Відповідь \\
\midrule
Обмежені ресурси & немає агента & мережевий моніторинг \\
Специфічні протоколи & сигнатури незастосовні & спеціалізовані правила \\
Складність оновлення & незакриті вразливості & сегментація, компенсація \\
Ціна збою & пріоритет доступності & захист без втручання \\
\bottomrule
\end{tabularx}
\end{center}

Спільним стратегічним принципом стає \term{ешелонований захист без покладання на кінцеву точку}: оскільки сам пристрій часто незахищуваний, основним рубежем стають сегментація мережі (ізоляція спеціалізованого сегмента), контроль і моніторинг трафіку на межах сегментів, гартування конфігурації й сувора профілактика. Там, де не можна виявити ЗПК \emph{на} пристрої, його виявляють \emph{навколо} пристрою --- за аномаліями мережевої взаємодії.

Ще один спільний принцип --- \emph{передбачуваність} спеціалізованих середовищ, яку захисник може обернути собі на користь. На відміну від універсальної робочої станції, де користувач запускає довільне ПЗ, спеціалізований пристрій зазвичай виконує вузький, сталий набір функцій: контролер циклічно опитує датчики й керує виконавчими механізмами, камера передає відеопотік на сервер, а не сканує сусідів. Ця вузькість означає, що нормальна поведінка добре окреслена, а отже будь-яке відхилення --- нетипове з'єднання, нова команда, сплеск трафіку --- підозріле з високою питомою вагою. Там, де на універсальній системі аномалія тоне в розмаїтті легітимної активності (проблема базової частоти, розділ~\ref{chap:ch03}), у спеціалізованому середовищі та сама аномалія різко виділяється. Тому \term{білий список дозволеної поведінки} (allow-listing) --- як мережевої, так і процесної --- тут значно дієвіший, ніж на настільних системах, і становить природний фундамент захисту.

\section{Промислові системи керування (ICS/SCADA)}\label{sec:ch14-ics}

\term{Промислові системи керування} (ICS, зокрема SCADA) керують фізичними процесами --- енергетикою, водопостачанням, виробництвом. Їх безпека~\cite{ch14:stouffer,ch14:iec62443} має принципову відмінність від IT: пріоритет --- \emph{доступність і цілісність} процесу (а не конфіденційність даних), бо збій означає фізичні наслідки. Специфіка: застарілі системи без вбудованої безпеки, спеціалізовані протоколи, неперервність роботи (неможливість зупинити для оновлення), конвергенція IT/OT (поєднання офісних і технологічних мереж створює нові вектори).

ЗПК для ICS --- вершина складності загроз. Stuxnet~\cite{ch14:langner} уперше показав код, що фізично руйнував обладнання (центрифуги), маніпулюючи логікою контролерів і приховуючи це від операторів. В Україні Industroyer/Industroyer2~\cite{ch14:industroyer} --- спеціалізований код для атак на енергосистему --- знав промислові протоколи й безпосередньо керував підстанціями, спричинивши вимкнення електропостачання. Ці кейси засвідчили: ЗПК для ICS потребує глибокого знання конкретного процесу й обладнання, діє цілеспрямовано (APT, розділ~\ref{chap:ch13}) і має на меті фізичний вплив.

Протидія: сувора сегментація IT/OT (демілітаризовані зони, односпрямовані шлюзи), моніторинг промислових протоколів на аномалії (нетипові команди контролерам), гартування й контроль цілісності логіки контролерів, фізичні запобіжники, незалежні від керуючого коду. Оскільки встановити агент на контролер зазвичай неможливо, ключовим є пасивний мережевий моніторинг OT-сегмента. Українські дослідження систематизують захист АСК ТП критичної інфраструктури~\cite{ch14:ics-ua}.

\begin{tikzfig}
\begin{tikzpicture}[font=\scriptsize,
  zn/.style={draw,rounded corners=2pt,minimum height=6mm,text width=22mm,align=center,font=\tiny},
  >={Stealth[length=1.4mm]}]
\node[zn] (it) at (0,0) {IT-мережа (офіс, інтернет)};
\node[zn,fill=black!6] (dmz) at (3.2,0) {ДМЗ / односпрямований шлюз};
\node[zn] (ot) at (6.4,0) {OT-мережа (SCADA)};
\node[zn] (plc) at (6.4,-1.3) {контролери (PLC)};
\node[zn] (phys) at (6.4,-2.6) {фізичні запобіжники};
\draw[<->] (it)--(dmz); \draw[->] (dmz)--(ot); \draw[->] (ot)--(plc); \draw[-] (plc)--(phys);
\node[font=\tiny,align=center] at (3.2,-1.3) {моніторинг\\протоколів};
\draw[->,dashed] (3.2,-1.0) -- (ot.south west);
\end{tikzpicture}
\figcap{Рис.~\thechapter.2 --- Сегментація IT/OT із демілітаризованою зоною, мережевим моніторингом і фізичними запобіжниками}
\end{tikzfig}

Специфіку ризику в ICS формалізує врахування \emph{фізичних} наслідків. На відміну від IT, де збиток вимірюють компрометацією даних, в OT очікуваний збиток зважують на тяжкість фізичного наслідку:

\begin{equation}
R_{\text{OT}} = \sum_{s} P(s)\cdot I_{\text{safety}}(s),
\label{eq:ch14-otrisk}
\end{equation}

де $P(s)$ --- імовірність сценарію $s$, а $I_{\text{safety}}(s)$ --- його наслідок із урахуванням загрози людям і довкіллю. Формула~\eqref{eq:ch14-otrisk} пояснює, чому в ICS навіть малоймовірні сценарії з катастрофічним $I_{\text{safety}}$ домінують у ризику й виправдовують дорогі незалежні фізичні запобіжники: жодне програмне виявлення не замінює захисту, що діє, навіть коли керуючий код повністю скомпрометовано.

\section{Інтернет речей (IoT)}\label{sec:ch14-iot}

\term{Інтернет речей} --- мільярди під'єднаних пристроїв (камери, маршрутизатори, датчики) --- став масовою ціллю через слабку безпеку: типові паролі, незакриті вразливості, відсутність оновлень, мінімальні ресурси для захисту~\cite{ch14:iot-owasp}. Ботнет Mirai~\cite{ch14:antonakakis} продемонстрував масштаб загрози: заразивши сотні тисяч пристроїв простим перебором типових паролів, він здійснив одні з найпотужніших DDoS-атак в історії. Привабливість IoT для зловмисника: масовість (величезна поверхня атаки), слабкий захист, цілодобова доступність, слабка видимість (власники не моніторять).

Швидкість захоплення IoT-ботнетом описують епідемічними моделями (розділ~\ref{chap:ch03}). На ранній стадії, коли вразлива популяція ще велика, кількість заражених пристроїв зростає майже експоненційно:

\begin{equation}
\frac{dN}{dt} = \beta N\left(1 - \frac{N}{N_{\max}}\right),
\label{eq:ch14-growth}
\end{equation}

де $N$ --- заражені пристрої, $N_{\max}$ --- розмір уразливої популяції, $\beta$ --- темп поширення (залежить від швидкості сканування й частки пристроїв із типовими паролями). Логістична форма~\eqref{eq:ch14-growth} пояснює вибуховий характер Mirai: за великого $N_{\max}$ (мільйони пристроїв із типовими паролями) і високого $\beta$ ботнет досягає масштабу, достатнього для рекордних DDoS, за години. Вона ж підказує важелі протидії: зниження $N_{\max}$ (зміна типових паролів, оновлення) і $\beta$ (обмеження сканування сегментацією) збивають епідемію ще на старті.

\begin{tikzfig}
\begin{tikzpicture}[font=\scriptsize,
  b/.style={draw,rounded corners=2pt,minimum height=5.5mm,text width=18mm,align=center,font=\tiny},
  >={Stealth[length=1.4mm]}]
\node[b] (scan) at (0,0) {сканування (типові паролі)};
\node[b] (inf) at (2.4,0) {зараження пристрою};
\node[b] (bot) at (4.8,0) {ботнет (C2)};
\node[b] (ddos) at (7.2,0) {DDoS / проксі / майнінг};
\draw[->] (scan)--(inf); \draw[->] (inf)--(bot); \draw[->] (bot)--(ddos);
\draw[->] (inf) to[out=-90,in=-90] node[below,font=\tiny]{поширення} (scan.south);
\end{tikzpicture}
\figcap{Рис.~\thechapter.1 --- Життєвий цикл IoT-ботнету на прикладі моделі Mirai}
\end{tikzfig}

Протидія ускладнена обмеженнями пристроїв: повноцінний захист на пристрої неможливий. Тому акцент --- на мережевому рівні (виявлення аномального трафіку, сегментація IoT в окремий сегмент, обмеження зовнішнього доступу), профілактиці (зміна типових паролів, оновлення, вимкнення зайвих служб) і безпеці за задумом (secure-by-design --- вимоги до виробників). Виявлення зараження спирається на аномалії поведінки (сплеск вихідного трафіку, звернення до командних серверів), а не на аналіз коду на пристрої. Вітчизняні роботи пропонують моделі виявлення шкідливої активності в IoT~\cite{ch14:iot-ua}.

\section{Мобільні платформи}\label{sec:ch14-mobile}

Мобільні ОС (Android, iOS) мають принципово іншу модель безпеки, ніж настільні: \term{пісочниця застосунків} (ізоляція кожного застосунку), дозвільна модель (контроль доступу до ресурсів), централізовані магазини з перевіркою, оновлення платформи~\cite{ch14:android}. Це ускладнює класичний ЗПК, але не унеможливлює: шкідливі застосунки обходять перевірку магазинів, зловживають дозволами, застосовують соціальну інженерію, експлуатують вразливості. Дослідження безпеки Android~\cite{ch14:enck} показали поширеність витоку даних і зловживання дозволами навіть у легітимних застосунках.

Специфіка мобільних загроз: шпигунський код (стеження, перехоплення повідомлень --- особливо небезпечний для цільових осіб), банківські трояни (накладні вікна для крадіжки облікових даних), зловживання дозволами. Протидія: модель безпеки платформи (пісочниця, дозволи), перевірка магазинів, мобільні засоби захисту, контроль у корпоративному середовищі (MDM), аналіз застосунків (статичний і динамічний, розділи~\ref{chap:ch06}--\ref{chap:ch07}, адаптовані до байт-коду й специфіки платформи). Ключова лінія --- дозвільна модель і пісочниця, що обмежують вплив навіть зловмисного застосунку. Порівняння моделей захисту чотирьох середовищ наведено в табл.~\thechapter.2.

\tabhead{Таблиця~\thechapter.2 --- Спеціалізовані середовища: модель захисту}

\begin{center}\footnotesize
\begin{tabularx}{\linewidth}{L{16mm} Y L{20mm}}
\toprule
Середовище & Основний рубіж & Пріоритет \\
\midrule
ICS/SCADA & сегментація, монітор OT & доступність, безпека \\
IoT & мережевий рівень, профілактика & масштаб, видимість \\
Мобільні & пісочниця, дозволи & ізоляція застосунків \\
Хмара/контейнери & цілісність образів, незмінність & динаміка, масштаб \\
\bottomrule
\end{tabularx}
\end{center}

\section{Хмарні та контейнерні середовища}\label{sec:ch14-cloud}

Хмарні~\cite{ch14:cloud-nist} й контейнерні~\cite{ch14:container} середовища зміщують модель загроз: спільна відповідальність (безпека розподілена між провайдером і клієнтом), ефемерність (контейнери живуть короткий час), масштаб і динамічність, нові вектори (компрометація образів, оркестрації, викрадення обчислень для майнінгу). ЗПК у хмарі: заражені образи контейнерів (ланцюг постачання, розділ~\ref{chap:ch15}), код для криптомайнінгу на викрадених ресурсах, компрометація керуючої площини.

Протидія адаптує знайомі принципи до динаміки хмари: контроль цілісності й сканування образів \emph{до} розгортання (а не лише під час виконання), найменші привілеї для робочих навантажень, сегментація й мікросегментація, моніторинг поведінки контейнерів (аномальні процеси, мережа), незмінна інфраструктура (immutable --- заміна замість зміни, що ускладнює персистентність). Ефемерність контейнерів --- палиця з двома кінцями: з одного боку, ускладнює персистентність ЗПК; з іншого, ускладнює розслідування (контейнер зникає разом зі слідами), що вимагає завчасного збору телеметрії.

Модель \emph{спільної відповідальності} заслуговує окремого наголосу, бо її нерозуміння --- типове джерело компрометацій. Провайдер відповідає за безпеку \emph{хмари} (фізичну інфраструктуру, гіпервізор, керуючі служби), клієнт --- за безпеку \emph{у хмарі} (конфігурацію, дані, застосунки, права доступу). Значна частина хмарних інцидентів спричинена не проломом захисту провайдера, а помилками клієнта: відкритими сховищами, надмірними правами, скомпрометованими ключами доступу, невиправленими застосунками. Для ЗПК це означає нові вектори, специфічні для хмари: викрадення хмарних облікових даних відкриває доступ до керуючої площини, звідки можна розгортати власні навантаження (майнінг, проксі) коштом жертви. Тому захист у хмарі зміщує акцент із «антивірусу» на керування ідентичністю, правами й конфігурацією, а виявлення --- на аномалії в журналах керуючої площини, а не лише в самих робочих навантаженнях.

\section{Кейс: захист промислового об'єкта критичної інфраструктури}\label{sec:ch14-case}

Розглянемо побудову захисту промислового об'єкта від цілеспрямованого ЗПК на кшталт Industroyer (захисна задача, за принципами IEC~62443~\cite{ch14:iec62443}).

\emph{Крок 1. Модель загроз.} Визначаємо активи (контролери, інженерні станції, керуючий процес), ймовірного супротивника (APT державного рівня, розділ~\ref{chap:ch13}), наслідки (фізичний збій, загроза людям) та вектори (конвергенція IT/OT, ланцюг постачання, знімні носії).

\emph{Крок 2. Сегментація.} Розділяємо IT- та OT-мережі демілітаризованою зоною й суворо контрольованими (ідеально --- односпрямованими) шлюзами, щоб компрометація офісної мережі не давала прямого доступу до OT.

\emph{Крок 3. Моніторинг OT.} Розгортаємо пасивний моніторинг промислових протоколів, що виявляє аномальні команди контролерам (нетипові для нормального процесу) --- ключовий рубіж там, де агент на контролері неможливий.

\emph{Крок 4. Гартування й цілісність.} Гартуємо конфігурацію інженерних станцій, контролюємо цілісність логіки контролерів, обмежуємо знімні носії й віддалений доступ, застосовуємо найменші привілеї.

\emph{Крок 5. Фізичні запобіжники й реагування.} Передбачаємо незалежні від керуючого коду фізичні запобіжники (щоб код не міг спричинити катастрофу) і план реагування (розділ~\ref{chap:ch16}) з урахуванням неперервності процесу.

Кейс ілюструє адаптацію загальних принципів до ICS: оскільки виявлення на контролері неможливе, захист будують навколо нього --- сегментацією, мережевим моніторингом, гартуванням і фізичними запобіжниками, з пріоритетом доступності й безпеки процесу.

\section*{Висновки до розділу}
\addcontentsline{toc}{section}{Висновки до розділу}

Спеціалізовані середовища --- ICS/SCADA, IoT, мобільні платформи, хмара --- мають власні моделі загроз і обмеження, через що прямий перенос «настільних» підходів часто неможливий (табл.~\thechapter.1). Спільний виклик --- неможливість повноцінного захисту на самому пристрої --- зміщує акцент на мережевий рівень, сегментацію й профілактику: там, де не можна виявити ЗПК на пристрої, його виявляють навколо нього. У ICS пріоритет доступності й фізичні наслідки роблять код на кшталт Stuxnet та Industroyer вершиною складності загроз; протидія --- сувора сегментація IT/OT, моніторинг промислових протоколів, гартування й фізичні запобіжники. IoT через масовість і слабкий захист став базою ботнетів (Mirai); протидія --- мережевий рівень, профілактика й безпека за задумом. Мобільні платформи мають сильну вбудовану модель безпеки (пісочниця, дозволи), що обмежує, але не усуває загрози. Хмарні й контейнерні середовища вимагають адаптації принципів до ефемерності й динаміки з акцентом на цілісність образів і незмінну інфраструктуру. Для України захист критичної інфраструктури є питанням національної безпеки. Головний висновок: протидія в спеціалізованих середовищах будується не на кінцевій точці, а навколо неї --- сегментацією, мережевим моніторингом аномалій, профілактикою й архітектурними запобіжниками, адаптованими до специфіки кожного середовища.

\selfcheck
\begin{enumerate}
\item Які спільні виклики відрізняють спеціалізовані середовища від класичних систем (табл.~\thechapter.1)?
\item Поясніть принцип ешелонованого захисту без покладання на кінцеву точку. Чому він стає спільним для спеціалізованих середовищ?
\item Чим безпека ICS/SCADA принципово відрізняється від IT? Чому пріоритетом є доступність і цілісність?
\item Опишіть кейси Stuxnet та Industroyer. Чому ЗПК для ICS є вершиною складності загроз?
\item Які заходи протидії застосовують у ICS? Чому ключовим є пасивний мережевий моніторинг OT?
\item Чому IoT став масовою ціллю? Опишіть життєвий цикл ботнету Mirai (рис.~\thechapter.1).
\item Як протидіють загрозам IoT за умов обмежених ресурсів пристроїв?
\item Чим модель безпеки мобільних платформ відрізняється від настільних? Як пісочниця й дозволи обмежують загрози?
\item Які специфічні мобільні загрози існують? Які заходи протидії застосовують?
\item Як хмарні й контейнерні середовища змінюють модель загроз? Що таке незмінна інфраструктура та чому вона ускладнює персистентність?
\item Чому ефемерність контейнерів є палицею з двома кінцями для захисту?
\item Опишіть побудову захисту промислового об'єкта (кейс підрозділу~\ref{sec:ch14-case}). Чому захист будують навколо контролера, а не на ньому?
\item Чому захист критичної інфраструктури є питанням національної безпеки України?
\item Дослідницьке: розробіть модель загроз і схему сегментації для обраного типу об'єкта критичної інфраструктури за IEC~62443 та оцініть залишковий ризик.
\item Дослідницьке: порівняйте ефективність мережевого виявлення аномалій для IoT-ботнету за різних наборів ознак трафіку на реалістичних даних.
\end{enumerate}
